% Strona tytułowa — uzupełnij poniższe pola.
%
% Tytuł: krótki, wraps w raggedright 9pt. Jeśli długi — rozważ skrót
% w nagłówku i pełen tytuł jako \section* na początku treści.
%
% Logo: umieść slayer-logo.png w katalogu wariantu.
% Formatka wykryje brak pliku i wstawi typograficzny placeholder.

\slayertitle
  {Tytuł dokumentu}%                    ← krótki tytuł
  {Autor · Kolektyw Slayer · v0.1}%     ← autor / wersja
  {Data}%                               ← np. Czerwiec 2026

% Abstract full-width (opcjonalny — usuń jeśli zbędny)
\noindent
{\fontsize{6}{6}\selectfont\ttfamily\color{ALERTRED!70}%
  [\,ABSTRACT\,·\,LANG:PL\,]}\par\vspace{2pt}
{\fontsize{8.5}{12}\selectfont\noindent
Krótkie streszczenie (3--5 zdań): co badamy, na jakich danych,
jaki wynik. Empiria, nie deklaracje.
\corrk\;{\fontsize{5.5}{5.5}\selectfont\ttfamily\color{BETON}STATUS:OK}}

\vspace{4pt}
{\color{sLightGrid}\hrule height 0.3pt}\vspace{1pt}%
{\color{ALERTRED}\hrule height 0.8pt}
\vspace{8pt}
